\documentclass[../main.tex]{subfiles}
\begin{document}

\section{S7comm detection rule design approach}
\label{sec:detection-overview}

This chapter defines a Suricata rule set~\cite{suricata_oisf,veit2021ics_dissectors} to detect S7comm attack behaviour documented in Chapter~\ref{sec:attack-scenarios}. Detection is signature-based: each rule tests fixed byte patterns in the TCP/102 payload spanning TPKT, COTP, the S7 header, and---where applicable---function codes or rate limits.

Rules are grouped by SID range: session setup and reconnaissance (\texttt{100000x}), read/write variables (\texttt{100001x}), block upload/download (\texttt{100002x}), Start/Stop control (\texttt{100003x}), DoS and rate floods (\texttt{100004x}), and malformed PDUs (\texttt{100005x}).

Each signature must specify three design choices: (i)~which encapsulation bytes are checked, (ii)~traffic direction relative to the PLC on port~102, and (iii)~whether a single match or a rate threshold is appropriate. The chapter follows that order: shared offsets, translation principles, then one subsection per Chapter~3 scenario.

\section{Translation principles}
\label{sec:pcap-to-rules}

Signatures in this chapter are derived from traffic observed in the Chapter~3 scenarios. Translation from a captured PDU to a Suricata rule is not a direct Wireshark export~\cite{wireshark_s7comm}; it applies the same logic to every scenario:

\begin{enumerate}
    \item \textbf{Identify protocol anchors.} Valid S7comm Data TPDUs begin with TPKT version \texttt{03 00} at offset~0, COTP Data TPDU \texttt{02 F0 80} at offset~4, and S7 protocol identifier \texttt{32} at offset~7. Missing anchors indicate non-S7 traffic or malformed samples (Section~\ref{sec:rules-malformed}).

    \item \textbf{Fix direction and function offset.} Client Jobs use \texttt{ROSCTR=0x01} (\texttt{|02 F0 80 32 01|} at offset~4); the function code lies at offset~17. Userdata (Read SZL, List blocks) uses \texttt{ROSCTR=0x07} with the same offset~17 layout. PLC Ack\_Data responses use \texttt{ROSCTR=0x03}; two error-class bytes precede the parameter block, so the function code moves to offset~19. Apply \texttt{flow:to\_server,established} or \texttt{flow:from\_server,established} accordingly before \texttt{content} matches.

    \item \textbf{Choose single-match vs.\ rate counting.} Isolated Setup or Write Var operations are normal engineering traffic and are detected with plain \texttt{content} rules. High-frequency repetition of valid PDUs indicates DoS and requires \texttt{threshold}. Malformed traffic alters one structural field; \texttt{content}, \texttt{byte\_test}, or \texttt{dsize} on that field suffices without a threshold.

    \item \textbf{Validate against recorded traffic.} Each signature is checked offline against scenario captures so alerts align with the attack phase. Missed matches usually trace to wrong offset (17 vs.\ 19), inverted flow direction, or incomplete visibility of the attack path on the monitored interface.
\end{enumerate}

The same function byte can therefore appear in multiple rules with different \texttt{flow} and \texttt{ROSCTR} tests---for example \texttt{0x1b} on opposite legs of the download handshake (Section~\ref{sec:down_up_program}).

\section{Common rules for identifying the S7comm protocol}
\label{sec:rule-foundation}

On TCP/102, S7comm rides inside TPKT and COTP before the S7 header~\cite{rfc905,rfc1006,rfc2126,snap7_s7protocol}. Table~\ref{tab:s7-payload-offsets} lists the payload offsets used throughout this chapter.

\begin{table}[H]
    \centering
    \caption{Byte offsets in the TCP/102 payload (S7comm Data TPDU)}
    \label{tab:s7-payload-offsets}
    \begin{tabular}{|c|c|p{6.8cm}|}
        \hline
        \textbf{Offset} & \textbf{Fixed or example value} & \textbf{Field} \\
        \hline
        0--1 & \texttt{03 00} & TPKT version and reserved \\
        \hline
        2--3 & \emph{varies} & TPKT length (big-endian; must match TCP payload for valid PDUs) \\
        \hline
        4--6 & \texttt{02 F0 80} & COTP Data TPDU \\
        \hline
        7 & \texttt{32} & S7 protocol identifier \\
        \hline
        8 & \texttt{01}/\texttt{03}/\texttt{07} & \texttt{ROSCTR}: Job / Ack\_Data / Userdata \\
        \hline
        17 & \texttt{0xF0}, \texttt{0x05}, \texttt{0x28}, \ldots & Function code on Job or Userdata PDUs (opcode depends on command) \\
        \hline
        19 & \texttt{0x04}, \texttt{0x1B}, \ldots & Function code on Ack\_Data PDUs (same role as offset~17, shifted by two error bytes) \\
        \hline
    \end{tabular}
\end{table}

Rows at offsets~17 and~19 are not single fixed bytes like the TPKT or COTP anchors: each S7comm command uses its own opcode at that position. Rules in this chapter match specific values observed in Chapter~3 (for example \texttt{0xF0} for Setup, \texttt{0x05} for Write Var). Bytes~2--3 encode TPKT length; malformed rules compare that field against the actual TCP segment size via \texttt{byte\_test} and \texttt{dsize}.

\subsection{Flow direction and PDU type conditions}
\label{sec:rule-flow-pdu}

In the lab, the PLC listens on TCP/102 and acts as the server; HMI and attacker hosts are clients. Before matching a function code, Suricata rules apply two filters: a \texttt{flow} keyword that fixes direction relative to the server port, and a \texttt{content} keyword that confirms PDU type and places subsequent matches at offset~17 or~19.

\paragraph{Job PDU --- client $\rightarrow$ PLC.}
Client operations---Setup Communication, Write Var, Request Download, PLC Control---use Job PDUs with \texttt{ROSCTR=0x01}. Matching rules combine:
\begin{itemize}
    \item \texttt{flow:to\_server,established}: packets sent to port~102 on \texttt{\$HOME\_NET} inside an established TCP session;
    \item \texttt{content:"|02 F0 80 32 01|"; offset:4; depth:5}: COTP Data TPDU \texttt{02 F0 80}, protocol id \texttt{32}, and Job \texttt{ROSCTR} \texttt{0x01}.
\end{itemize}
Userdata commands (Read SZL, List blocks) set \texttt{ROSCTR=0x07}; the header pattern becomes \texttt{|02 F0 80 32 07|} at offset~4, with the function byte still at offset~17.

\paragraph{Ack\_Data PDU --- PLC $\rightarrow$ client.}
PLC responses---write acknowledgements, uploaded block data---use \texttt{ROSCTR=0x03}. Matching rules combine:
\begin{itemize}
    \item \texttt{flow:from\_server,established}: packets sourced from the PLC on port~102;
    \item \texttt{content:"|02 F0 80 32 03|"; offset:4; depth:5}: Ack\_Data header confirmation.
\end{itemize}
The Ack\_Data parameter block follows two error-class bytes, shifting the function code from offset~17 to offset~19. Response-direction rules (\texttt{sid:1000026}, \texttt{sid:1000013}) rely on this offset consistently.

These filters discard non-S7 traffic on port~102, separate requests from responses, and anchor \texttt{content} before scenario-specific signatures in Section~\ref{sec:rules-by-scenario}.

\section{Rules by attack scenario}
\label{sec:rules-by-scenario}

\subsection{Reconnaissance}
\label{sec:rules-recon}

Based on Section~\ref{sec:reconnaissance}: TCP/102 scanning, Read SZL (\texttt{0x0011}, \texttt{0x001c}), and block listing.

\paragraph{Detection rationale.}
Reconnaissance traffic begins with Setup Communication (\texttt{0xF0} at offset~17 on a Job PDU). Read SZL requests follow as Userdata (\texttt{ROSCTR=0x07}) with parameter head \texttt{00 01 12 04 11 44 01 00} at offset~17. Module identification reads SZL-ID \texttt{0x0011} (\texttt{00 11 00 01}); communication status uses \texttt{0x001c} (\texttt{00 1C 00 00}). Block inventory uses subfunction List blocks with parameter block \texttt{00 01 12 04 11 13 01 00} instead of \texttt{...44...}. These byte strings are stable across S7 stacks and therefore become \texttt{content} anchors rather than IP-based rules.

\paragraph{Rule \texttt{sid:1000001}.}
\begin{packetdump}
alert tcp any any -> $HOME_NET 102 (msg:"ICS S7COMM Setup Communication to PLC";
  flow:to_server,established; content:"|03 00|"; offset:0; depth:2;
  content:"|02 F0 80 32 01|"; offset:4; depth:5; content:"|F0|";
  offset:17; depth:1; sid:1000001; rev:1;
  metadata:service iso-tsap, protocol s7comm, attack reconnaissance; priority:3;)
\end{packetdump}

\paragraph{Rule design.}
The first S7comm Job after COTP setup carries \texttt{0xF0} at offset~17 while bytes 4--8 read \texttt{02 F0 80 32 01}. That byte is the Setup Communication function code, required before any SZL or block query. The rule chains TPKT, Job header, and \texttt{|F0|} with \texttt{flow:to\_server} because only the client sends this PDU to the PLC. It fires on every new session setup, including legitimate HMI connections, so it is a context alert (priority~3) rather than a standalone attack indicator.

\paragraph{Rule \texttt{sid:1000002}.}
\begin{packetdump}
alert tcp any any -> $HOME_NET 102 (msg:"ICS S7COMM Read SZL request - PLC
  information enumeration"; flow:to_server,established; content:"|03 00|";
  offset:0; depth:2; content:"|02 F0 80 32 07|"; offset:4; depth:5;
  content:"|00 01 12 04 11 44 01 00|"; offset:17; depth:8; sid:1000002; rev:1;
  metadata:service iso-tsap, protocol s7comm, attack reconnaissance; priority:2;)
\end{packetdump}

\paragraph{Rule design.}
Read SZL appears as Userdata with \texttt{ROSCTR=0x07} and an eight-byte parameter prefix \texttt{00 01 12 04 11 44 01 00} (method \texttt{0x11}, function group CPU functions, subfunction Read SZL). Matching the full eight-byte block at offset~17 avoids confusing Read SZL with other Userdata services that share \texttt{32 07} but use a different subfunction byte. The rule does not distinguish which SZL-ID is requested; that refinement is left to \texttt{sid:1000003} and \texttt{sid:1000004}.

\paragraph{Rule \texttt{sid:1000003}.}
\begin{packetdump}
alert tcp any any -> $HOME_NET 102 (msg:"ICS S7COMM Read SZL module identification
  0x0011"; flow:to_server,established; content:"|03 00|"; offset:0; depth:2;
  content:"|02 F0 80 32 07|"; offset:4; depth:5;
  content:"|00 01 12 04 11 44 01 00|"; offset:17; depth:8;
  content:"|00 11 00 01|"; distance:0; within:12; sid:1000003; rev:1;
  metadata:service iso-tsap, protocol s7comm, attack reconnaissance; priority:2;)
\end{packetdump}

\paragraph{Rule design.}
The data section after the Read SZL parameter contains SZL-ID \texttt{0x0011} and index \texttt{0x0001}, encoded as \texttt{|00 11 00 01|} within twelve bytes of the parameter head. That read returns the module order number used in device fingerprinting. The rule adds this four-byte constraint with \texttt{distance:0; within:12} so generic Read SZL probes (\texttt{sid:1000002}) and module-fingerprint reads are separated. Engineering tools may also read \texttt{0x0011} during commissioning, so the alert should be correlated with source address and timing.

\paragraph{Rule \texttt{sid:1000004}.}
\begin{packetdump}
alert tcp any any -> $HOME_NET 102 (msg:"ICS S7COMM Read SZL communication status
  0x001c"; flow:to_server,established; content:"|03 00|"; offset:0; depth:2;
  content:"|02 F0 80 32 07|"; offset:4; depth:5;
  content:"|00 01 12 04 11 44 01 00|"; offset:17; depth:8;
  content:"|00 1C 00 00|"; distance:0; within:12; sid:1000004; rev:1;
  metadata:service iso-tsap, protocol s7comm, attack reconnaissance; priority:2;)
\end{packetdump}

\paragraph{Rule design.}
A follow-on probe requests SZL-ID \texttt{0x001c} (communication status partial list), encoded as \texttt{|00 1C 00 00|} after the Read SZL parameter block. The PLC reply supplies serial number, system name, and module type strings used in device enumeration. This rule targets that specific step rather than every SZL read. It will not fire on \texttt{0x0011} reads matched by \texttt{sid:1000003}, but any tool that reads \texttt{0x001c} for diagnostics will also match.

\paragraph{Rule \texttt{sid:1000005}.}
\begin{packetdump}
alert tcp any any -> $HOME_NET 102 (msg:"ICS S7COMM List Blocks request - PLC block
  inventory enumeration"; flow:to_server,established; content:"|03 00|";
  offset:0; depth:2; content:"|02 F0 80 32 07|"; offset:4; depth:5;
  content:"|00 01 12 04 11 13 01 00|"; offset:17; depth:8; sid:1000005; rev:1;
  metadata:service iso-tsap, protocol s7comm, attack reconnaissance; priority:2;)
\end{packetdump}

\paragraph{Rule design.}
Block inventory sends Userdata with function group Block functions and subfunction List blocks; the eight-byte parameter at offset~17 is \texttt{|00 01 12 04 11 13 01 00|}, differing from Read SZL only in the subfunction byte (\texttt{0x13} vs \texttt{0x44}). An attacker mapping OB/DB/FC/FB before upload or download will trigger this rule after Setup Communication. TIA Portal block browsing uses the same S7comm service, so the alert alone does not prove malicious intent.

\subsection{Start/Stop PLC command replay}
\label{sec:rules-start-stop}

Based on Section~\ref{sec:start_stop_plc}: replay of \texttt{PLC Control 0x28} and \texttt{PLC Stop 0x29} payloads with PI-Service \texttt{P\_PROGRAM}.

\paragraph{Detection rationale.}
TIA Portal Start CPU and Stop CPU payloads from Section~\ref{sec:start_stop_plc} begin with \texttt{0300...02f0803201} (TPKT + COTP + Job header). At offset~17, Start uses function \texttt{0x28} and Stop uses \texttt{0x29}. The ASCII string \texttt{P\_PROGRAM} (\texttt{50 5F 50 52 4F 47 52 41 4D} in hex) appears in the PI-Service data within 24 bytes after the function byte. Rules \texttt{sid:1000030}/\texttt{sid:1000032} match only the function byte; \texttt{sid:1000031}/\texttt{sid:1000033} add the \texttt{P\_PROGRAM} string to match replayed payloads exactly.

\paragraph{Rule \texttt{sid:1000030}.}
\begin{packetdump}
alert tcp any any -> $HOME_NET 102 (msg:"ICS S7COMM PLC Control command";
  flow:to_server,established; content:"|03 00|"; offset:0; depth:2;
  content:"|02 F0 80 32 01|"; offset:4; depth:5; content:"|28|";
  offset:17; depth:1; sid:1000030; rev:1;
  metadata:service iso-tsap, protocol s7comm, attack plc_control; priority:1;)
\end{packetdump}

\paragraph{Rule design.}
The captured Start CPU payload places \texttt{0x28} (PLC Control) at offset~17 on a Job PDU directed to port~102. The same function code is used for several PI-Services besides cold start---activate block, copy RAM to ROM, and others---so the rule intentionally omits parameter bytes. It gives early warning whenever any Control command reaches the PLC from an unexpected client. Legitimate engineering stations also send \texttt{0x28} during commissioning, which limits specificity compared to \texttt{sid:1000031}.

\paragraph{Rule \texttt{sid:1000031}.}
\begin{packetdump}
alert tcp any any -> $HOME_NET 102 (msg:"ICS S7COMM replayed Start CPU P_PROGRAM";
  flow:to_server,established; content:"|03 00|"; offset:0; depth:2;
  content:"|02 F0 80 32 01|"; offset:4; depth:5; content:"|28|";
  offset:17; depth:1; content:"P_PROGRAM"; distance:0; within:24;
  sid:1000031; rev:1;
  metadata:service iso-tsap, protocol s7comm, attack start_stop_replay; priority:1;)
\end{packetdump}

\paragraph{Rule design.}
The Start replay payload in Section~\ref{sec:start_stop_plc} contains \texttt{0x28} followed within 24 bytes by the ASCII PI-Service name \texttt{P\_PROGRAM}, which selects the CPU program for start. The rule inherits the \texttt{0x28} match from \texttt{sid:1000030} and adds a string match for the replayed payload. Other Control commands that use a different PI-Service name will not match, which reduces false positives from benign \texttt{0x28} traffic.

\paragraph{Rule \texttt{sid:1000032}.}
\begin{packetdump}
alert tcp any any -> $HOME_NET 102 (msg:"ICS S7COMM PLC Stop command";
  flow:to_server,established; content:"|03 00|"; offset:0; depth:2;
  content:"|02 F0 80 32 01|"; offset:4; depth:5; content:"|29|";
  offset:17; depth:1; sid:1000032; rev:1;
  metadata:service iso-tsap, protocol s7comm, attack plc_stop; priority:1;)
\end{packetdump}

\paragraph{Rule design.}
Stop CPU replay uses function \texttt{0x29} (PLC Stop) at offset~17 on the same Job header pattern \texttt{02 F0 80 32 01}. In the Stop hex dump (\texttt{0300002102f080...29...}), byte~17 is \texttt{0x29} regardless of PI-Service details. This rule alerts on any Stop request from client to PLC before checking whether the payload matches the recorded TIA Portal replay. Emergency stop from an authorized engineering station will also trigger it.

\paragraph{Rule \texttt{sid:1000033}.}
\begin{packetdump}
alert tcp any any -> $HOME_NET 102 (msg:"ICS S7COMM replayed Stop CPU P_PROGRAM";
  flow:to_server,established; content:"|03 00|"; offset:0; depth:2;
  content:"|02 F0 80 32 01|"; offset:4; depth:5; content:"|29|";
  offset:17; depth:1; content:"P_PROGRAM"; distance:0; within:24;
  sid:1000033; rev:1;
  metadata:service iso-tsap, protocol s7comm, attack start_stop_replay; priority:1;)
\end{packetdump}

\paragraph{Rule design.}
The Stop replay payload in Section~\ref{sec:start_stop_plc} pairs \texttt{0x29} with \texttt{P\_PROGRAM} in the same relative position as the Start packet. Matching both narrows the alert to replayed Stop traffic rather than any Stop command. If an attacker replays a Stop PDU without the \texttt{P\_PROGRAM} string, only \texttt{sid:1000032} fires.

\subsection{Block upload and download}
\label{sec:rules-up-download}

Based on Section~\ref{sec:down_up_program}: function sequence \texttt{0x1a}--\texttt{0x1f} on Siemens block-transfer traffic.

\paragraph{Detection rationale.}
Upload and download sessions in Section~\ref{sec:down_up_program} follow the Siemens function sequence: Request Download \texttt{0x1a}, Download Block \texttt{0x1b}, Download Ended \texttt{0x1c}, Start Upload \texttt{0x1d}, Upload Block \texttt{0x1e}, End Upload \texttt{0x1f}. Each opcode appears at offset~17 on Job PDUs from client to PLC, except when the PLC returns block data---then the client sends Ack\_Data (\texttt{ROSCTR=0x03}) and \texttt{0x1b} moves to offset~19 (\texttt{sid:1000026}). The PLC may also initiate Download Block (\texttt{0x1b}, Job, \texttt{from\_server}) as in \texttt{sid:1000021}. Live block-transfer PCAPs were not obtained in the lab; these rules were tuned on reference Wireshark captures~\cite{wireshark_s7comm} and documented hex sequences.

\paragraph{Rule \texttt{sid:1000020}.}
\begin{packetdump}
alert tcp any any -> $HOME_NET 102 (msg:"ICS S7COMM Request Download block to PLC";
  flow:to_server,established; content:"|03 00|"; offset:0; depth:2;
  content:"|02 F0 80 32 01|"; offset:4; depth:5; content:"|1A|";
  offset:17; depth:1; sid:1000020; rev:1;
  metadata:service iso-tsap, protocol s7comm, attack program_download; priority:1;)
\end{packetdump}

\paragraph{Rule design.}
The first Job in a download sequence carries function \texttt{0x1a} (Request Download) at offset~17 after the standard \texttt{02 F0 80 32 01} header. The rule uses \texttt{flow:to\_server} because only the engineering station sends Request Download. Normal program loads from TIA Portal begin with the same byte, so source IP and session context matter for interpretation.

\paragraph{Rule \texttt{sid:1000021}.}
\begin{packetdump}
alert tcp $HOME_NET 102 -> any any (msg:"ICS S7COMM PLC requests Download Block data";
  flow:from_server,established; content:"|03 00|"; offset:0; depth:2;
  content:"|02 F0 80 32 01|"; offset:4; depth:5; content:"|1B|";
  offset:17; depth:1; sid:1000021; rev:1;
  metadata:service iso-tsap, protocol s7comm, attack program_download; priority:1;)
\end{packetdump}

\paragraph{Rule design.}
During download, the PLC sometimes sends a Job PDU with \texttt{0x1b} (Download Block request) from port~102 toward the client---the opposite direction of \texttt{sid:1000026}. The rule reverses source and destination (\texttt{\$HOME\_NET 102 -> any any}) and uses \texttt{flow:from\_server} with \texttt{ROSCTR=0x01} and offset~17. Without the direction constraint, the same \texttt{0x1b} byte would match both handshake legs.

\paragraph{Rule \texttt{sid:1000026}.}
\begin{packetdump}
alert tcp any any -> $HOME_NET 102 (msg:"ICS S7COMM client sends Download Block data
  to PLC"; flow:to_server,established; content:"|03 00|"; offset:0; depth:2;
  content:"|02 F0 80 32 03|"; offset:4; depth:5; content:"|1B|";
  offset:19; depth:1; sid:1000026; rev:1;
  metadata:service iso-tsap, protocol s7comm, attack program_download; priority:1;)
\end{packetdump}

\paragraph{Rule design.}
When the client pushes block bytes to the PLC, the PDU type switches to Ack\_Data (\texttt{32 03}) and the function \texttt{0x1b} shifts to offset~19 because of the two-byte error header. This packet carries the actual program content in the data section. The pair \texttt{sid:1000021}/\texttt{sid:1000026} mirrors the two directions observed in the download sequence diagram of Section~\ref{sec:down_up_program}. The rule does not inspect block content---only the presence of a Download Block transfer.

\paragraph{Rule \texttt{sid:1000022}.}
\begin{packetdump}
alert tcp any any -> $HOME_NET 102 (msg:"ICS S7COMM Download Ended";
  flow:to_server,established; content:"|03 00|"; offset:0; depth:2;
  content:"|02 F0 80 32 01|"; offset:4; depth:5; content:"|1C|";
  offset:17; depth:1; sid:1000022; rev:1;
  metadata:service iso-tsap, protocol s7comm, attack program_download; priority:1;)
\end{packetdump}

\paragraph{Rule design.}
Download Ended (\texttt{0x1c}) closes the download session on a Job PDU from client to PLC. Matching \texttt{|1C|} at offset~17 marks completion of a program push. A single \texttt{sid:1000020} without follow-up \texttt{sid:1000026}/\texttt{sid:1000022} may indicate an aborted download rather than a full compromise.

\paragraph{Rule \texttt{sid:1000023}.}
\begin{packetdump}
alert tcp any any -> $HOME_NET 102 (msg:"ICS S7COMM Start Upload block from PLC";
  flow:to_server,established; content:"|03 00|"; offset:0; depth:2;
  content:"|02 F0 80 32 01|"; offset:4; depth:5; content:"|1D|";
  offset:17; depth:1; sid:1000023; rev:1;
  metadata:service iso-tsap, protocol s7comm, attack program_upload; priority:1;)
\end{packetdump}

\paragraph{Rule design.}
Start Upload (\texttt{0x1d}) opens an upload session: the client requests a block name such as \texttt{\_0800001P} from the PLC. The function byte at offset~17 is \texttt{0x1d} on a Job PDU. This is the first step when an attacker extracts logic from the device. Authorized backup operations from TIA Portal send the same PDU, so the alert indicates exfiltration only when the source is not the known engineering station.

\paragraph{Rule \texttt{sid:1000024}.}
\begin{packetdump}
alert tcp any any -> $HOME_NET 102 (msg:"ICS S7COMM Upload Block request";
  flow:to_server,established; content:"|03 00|"; offset:0; depth:2;
  content:"|02 F0 80 32 01|"; offset:4; depth:5; content:"|1E|";
  offset:17; depth:1; sid:1000024; rev:1;
  metadata:service iso-tsap, protocol s7comm, attack program_upload; priority:1;)
\end{packetdump}

\paragraph{Rule design.}
Upload Block (\texttt{0x1e}) repeats while the client pulls data chunks from the PLC; each request is a Job with \texttt{0x1e} at offset~17. The rule counts each request separately, so a full upload generates many alerts. It does not inspect block content returned by the PLC.

\paragraph{Rule \texttt{sid:1000025}.}
\begin{packetdump}
alert tcp any any -> $HOME_NET 102 (msg:"ICS S7COMM End Upload";
  flow:to_server,established; content:"|03 00|"; offset:0; depth:2;
  content:"|02 F0 80 32 01|"; offset:4; depth:5; content:"|1F|";
  offset:17; depth:1; sid:1000025; rev:1;
  metadata:service iso-tsap, protocol s7comm, attack program_upload; priority:1;)
\end{packetdump}

\paragraph{Rule design.}
End Upload (\texttt{0x1f}) terminates the upload session on a Job PDU. It typically follows a burst of \texttt{sid:1000024} alerts in the same TCP session. Together, \texttt{sid:1000023}--\texttt{sid:1000025} outline the upload timeline visible in Section~\ref{sec:down_up_program}. The rule alone does not reveal which block was copied.

\subsection{Process variable writes and MITM observation}
\label{sec:rules-stuxnet}

Based on Section~\ref{sec:stuxnet_mitm_sim}: process writes via Write Var (for example \texttt{DB1.DBW2}, \texttt{iSetpoint} \texttt{0x04b0}$\rightarrow$\texttt{0x07d0}); Ack\_Data interception toward the HMI. Rule \texttt{sid:1000011} matches any write to the DB memory area without fixing a DB number.

\paragraph{Detection rationale.}
The MITM scenario in Section~\ref{sec:stuxnet_mitm_sim} sends Write Var (\texttt{0x05} at offset~17) to change a datablock setpoint. The Job PDU contains request items with Transport size \texttt{0x12}, length \texttt{0x10}, and area byte \texttt{0x84} (DB area)---the pattern \texttt{|12|...|10|...|84|} used in \texttt{sid:1000011}. The MITM leg does not alter request opcodes; the interceptor rewrites Read Var responses on the PLC$\rightarrow$HMI path, increasing the rate of Ack\_Data (\texttt{ROSCTR=0x03}) without a single suspicious function byte. That pattern motivates \texttt{sid:1000013} with a threshold rather than a content match on modified payload bytes.

\paragraph{Rule \texttt{sid:1000010}.}
\begin{packetdump}
alert tcp any any -> $HOME_NET 102 (msg:"ICS S7COMM Write Var request to PLC";
  flow:to_server,established; content:"|03 00|"; offset:0; depth:2;
  content:"|02 F0 80 32 01|"; offset:4; depth:5; content:"|05|";
  offset:17; depth:1; sid:1000010; rev:1;
  metadata:service iso-tsap, protocol s7comm, attack modify_process_variable;
  priority:1;)
\end{packetdump}

\paragraph{Rule design.}
Every Write Var in the MITM scenario places function \texttt{0x05} at offset~17 on a Job PDU to port~102. The rule matches only the function byte, so it covers writes to M, I, Q, and DB alike. It is the broad base layer beneath the DB-specific \texttt{sid:1000011} and will fire on normal HMI setpoint updates unless filtered by source IP.

\paragraph{Rule \texttt{sid:1000011}.}
\begin{packetdump}
alert tcp any any -> $HOME_NET 102 (msg:"ICS S7COMM Write Var to DB memory area";
  flow:to_server,established; content:"|03 00|"; offset:0; depth:2;
  content:"|02 F0 80 32 01|"; offset:4; depth:5; content:"|05|";
  offset:17; depth:1; content:"|12|"; offset:19; depth:1;
  content:"|10|"; offset:21; depth:1; content:"|84|";
  distance:5; within:1; sid:1000011; rev:2;
  metadata:service iso-tsap, protocol s7comm, attack modify_process_variable;
  priority:1;)
\end{packetdump}

\paragraph{Rule design.}
The Write Var PDU for a DB address contains an Any-type request item with Transport size \texttt{0x12}, length \texttt{0x10}, and area code \texttt{0x84} (S7 DB area) five bytes after the Syntax ID field. The rule checks \texttt{|12|}, \texttt{|10|}, and \texttt{|84|} at offsets 19--21 plus \texttt{distance:5; within:1} for the area byte without hard-coding DB number or offset. Writes to M/I/Q using area \texttt{0x83}/\texttt{0x81}/\texttt{0x82} do not match. Legitimate DB tuning and attack writes still require contextual correlation.

\paragraph{Rule \texttt{sid:1000013}.}
\begin{packetdump}
alert tcp any 102 -> $HOME_NET any (msg:"ICS S7COMM suspicious Ack_Data response
  to HMI/engineering station"; flow:from_server,established;
  content:"|03 00|"; offset:0; depth:2; content:"|02 F0 80 32 03|";
  offset:4; depth:5; sid:1000013; rev:1;
  metadata:service iso-tsap, protocol s7comm, attack possible_mitm_observation;
  threshold:type both, track by_src, count 60, seconds 10; priority:3;)
\end{packetdump}

\paragraph{Rule design.}
During MITM (Section~\ref{sec:stuxnet_mitm_sim}), the HMI still polls the PLC and receives Ack\_Data PDUs (\texttt{02 F0 80 32 03}) from port~102; the interceptor may rewrite payload bytes, but the header pattern stays valid. Counting more than 60 such responses in ten seconds from one PLC source flags unusually heavy read traffic toward the HMI subnet. The rule does not inspect modified values and cannot prove payload tampering---it supports correlation with \texttt{sid:1000011} and L2/ARP alerts outside the S7comm rule set.

\subsection{Denial of service on the S7comm channel}
\label{sec:rules-dos}

Based on Section~\ref{sec:dos_s7comm}: repeated Setup Communication, TCP connect churn, and Read SZL at high rate. Unlike single-opcode rules, the \texttt{100004x} group uses \texttt{threshold} to alert when one source exceeds a limit within ten seconds.

\paragraph{Detection rationale.}
DoS traffic in Section~\ref{sec:dos_s7comm} contains three distinguishable phases. Phase~1 repeats Setup PDUs (\texttt{0xF0} at offset~17, often on new TCP sessions) at high rate---each individual packet is valid. Phase~2 produces many TCP SYN packets to port~102 without completing S7comm. Phase~3 reuses the Read SZL byte string from \texttt{sid:1000002} at 25+ times per ten seconds. A generic Job flood mixes multiple opcodes under \texttt{32 01}. Thresholds were chosen so legitimate HMI traffic stays below the alert counts defined in the rules below.

\paragraph{Rule \texttt{sid:1000040}.}
\begin{packetdump}
alert tcp any any -> $HOME_NET 102 (msg:"ICS possible TCP SYN flood to ISO-TSAP port 102";
  flow:to_server; flags:S;
  threshold:type both, track by_src, count 40, seconds 10; sid:1000040; rev:1;
  metadata:attack dos_connection_flood; priority:1;)
\end{packetdump}

\paragraph{Rule design.}
The TCP connect-churn phase opens and closes sessions rapidly, producing bursts of SYN packets to port~102 without a full S7comm handshake. This rule counts SYN segments to \texttt{\$HOME\_NET:102} per source and fires above 40 in ten seconds. It does not inspect S7 payload at all. Fast port scans or misconfigured clients can produce similar SYN rates, so the alert needs correlation with S7 flood phases on the same source IP.

\paragraph{Rule \texttt{sid:1000041}.}
\begin{packetdump}
alert tcp any any -> $HOME_NET 102 (msg:"ICS S7COMM possible Setup Communication flood to PLC";
  flow:to_server,established; content:"|03 00|"; offset:0; depth:2;
  content:"|02 F0 80 32 01|"; offset:4; depth:5; content:"|F0|"; offset:17; depth:1;
  threshold:type both, track by_src, count 25, seconds 10; sid:1000041; rev:1;
  metadata:attack dos_s7comm_setup_flood; priority:1;)
\end{packetdump}

\paragraph{Rule design.}
Each Setup flood packet is byte-identical to a normal Setup Communication---the same \texttt{|F0|} at offset~17 matched by \texttt{sid:1000001}. A single packet cannot distinguish flood from a legitimate setup; the threshold requires 25 or more Setup PDUs from one source in ten seconds. Reconnect loops on a faulty client could approach the threshold and need baseline tuning.

\paragraph{Rule \texttt{sid:1000042}.}
\begin{packetdump}
alert tcp any any -> $HOME_NET 102 (msg:"ICS S7COMM possible Read SZL flood - PLC service load";
  flow:to_server,established; content:"|03 00|"; offset:0; depth:2;
  content:"|02 F0 80 32 07|"; offset:4; depth:5;
  content:"|00 01 12 04 11 44 01 00|"; offset:17; depth:8;
  threshold:type both, track by_src, count 25, seconds 10; sid:1000042; rev:1;
  metadata:attack dos_s7comm_szl_flood; priority:1;)
\end{packetdump}

\paragraph{Rule design.}
The Read SZL flood phase sends the same Userdata header as reconnaissance; only the rate differs. At 25+ requests per ten seconds the PLC CPU spends measurable time building SZL lists. The content match is copied from \texttt{sid:1000002}; the threshold distinguishes attack from a single enumeration run. Aggressive monitoring tools that poll SZL periodically could false-positive if their interval is misconfigured.

\paragraph{Rule \texttt{sid:1000043}.}
\begin{packetdump}
alert tcp any any -> $HOME_NET 102 (msg:"ICS S7COMM possible S7 Job request flood to PLC";
  flow:to_server,established; content:"|03 00|"; offset:0; depth:2;
  content:"|02 F0 80 32 01|"; offset:4; depth:5;
  threshold:type both, track by_src, count 80, seconds 10; sid:1000043; rev:1;
  metadata:attack dos_s7comm_job_flood; priority:1;)
\end{packetdump}

\paragraph{Rule design.}
When mixed Job opcodes appear under \texttt{ROSCTR=0x01}, no single function byte is stable enough for a narrow signature. This rule counts 80 Job-header PDUs per ten seconds without checking offset~17, covering Setup floods and generic Job spam. The higher count avoids overlap with \texttt{sid:1000041} on Setup-only floods but may miss slow-rate abuse below 80 packets.

\subsection{Malformed packets and abnormal PDUs}
\label{sec:rules-malformed}

Malformed rules detect abnormal payload structure on TCP/102, not dangerous opcodes alone. Each fault type in Table~\ref{tab:malformed-changes} maps to a signature using \texttt{content}, \texttt{byte\_test}, or \texttt{dsize} at the offsets in Table~\ref{tab:s7-payload-offsets}; the subsections below list the corresponding SIDs.

\paragraph{Detection rationale.}
Each lab fault is derived from a valid Setup Communication template. Invalid TPKT version shows \texttt{02 00} instead of \texttt{03 00} at offset~0; length mismatch modes declare \texttt{01 00} or \texttt{00 08} while the TCP segment size differs; invalid COTP replaces \texttt{02 F0 80} with \texttt{01 00 00}; header faults change byte~7 to \texttt{0x31} or byte~8 to \texttt{0xFF}; function-layer faults set offset~17 to \texttt{0xFF} or send a truncated Write Var PDU; truncation sends only 14 bytes so offset~17 never arrives.

\paragraph{Group 1 --- TPKT / COTP (\texttt{sid:1000050}, \texttt{1000051}, \texttt{1000058}, \texttt{1000056}).}

\texttt{1000050}: \texttt{content:!|03 00|} at offset~0---catches non-standard TPKT version (\texttt{0x02 00}). \texttt{1000051}: after confirming \texttt{03 00}, \texttt{byte\_test} compares TPKT length (big-endian, offset~2) against \texttt{dsize}---fires when length is larger than the TCP payload (\texttt{0x0100}). \texttt{1000058}: same idea but length is smaller than the actual payload (\texttt{0x0008} with a full Setup body). \texttt{1000056}: valid TPKT but missing COTP \texttt{02 F0 80} at offset~4.

\paragraph{Rule design.}
Wrong TPKT version breaks dissection at the outer encapsulation; inflated or deflated length fields disagree with the TCP segment size; missing COTP prevents S7comm parsing altogether. These rules use negated content, \texttt{byte\_test}, or absence of \texttt{02 F0 80} rather than function opcodes because malformed packets may never reach offset~17. They can also match non-S7 traffic on port~102 if \texttt{HOME\_NET} is wide.

\paragraph{Group 2 --- S7 header (\texttt{sid:1000052}, \texttt{1000053}).}

\texttt{1000052}: COTP DT is present but Protocol ID byte at offset~7 is not \texttt{0x32}. \texttt{1000053}: matches \texttt{|02 F0 80 32 FF|} at offset~4 (\texttt{ROSCTR} \texttt{0xFF}); additional invalid \texttt{ROSCTR} values can be added with separate rules if needed.

\paragraph{Rule design.}
Once TPKT and COTP look valid, bytes~7--8 still identify S7comm. Invalid protocol ID uses \texttt{0x31} instead of \texttt{0x32} at offset~7; invalid \texttt{ROSCTR} sets byte~8 to \texttt{0xFF}. Matching these fixed patterns catches header fuzzing before any function code is evaluated.

\paragraph{Group 3 --- Function / parameters (\texttt{sid:1000054}, \texttt{1000057}).}

\texttt{1000054}: Job (\texttt{32 01}) with function \texttt{0xFF} at offset~17---opcode outside the lab command set. \texttt{1000057}: still a Job with \texttt{0x05} (Write Var) but \texttt{dsize<25}: PDU too short for a valid Write Var.

\paragraph{Rule design.}
Group~3 packets pass header checks but fail semantic validation on the PLC. Unknown function \texttt{0xFF} at offset~17 or a truncated Write Var PDU (\texttt{0x05} with total length below 25 bytes) triggers these rules.

\paragraph{Group 4 --- Truncation (\texttt{sid:1000055}).}

\texttt{1000055}: header \texttt{02 F0 80 32 01} is present but \texttt{dsize<18}---packet too short to contain the function byte at offset~17 (Setup truncated to 14 bytes). This is a ``missing byte'' signature independent of the specific opcode.

\paragraph{Rule design.}
Truncation sends only the first 14 bytes of a Setup PDU, leaving no byte at offset~17. \texttt{dsize<18} detects any Job header without a reachable function field, including random truncation by a proxy.

\paragraph{General note.}
Malformed rules do not replace DoS rules (\texttt{100004x}): flooding valid PDUs still requires \texttt{threshold}. Conversely, some malformed packets will not trigger \texttt{1000001} (Setup) because \texttt{0xF0} is missing or offset~17 is unreachable. In the lab evaluation (Chapter~\ref{sec:eval-malformed}), a malformed scenario is considered detected when at least one \texttt{100005x} alert appears, without requiring the PLC to change RUN/STOP state.

\end{document}
